% =============================================================
%  PDL — Topological Origin of the Exponent 18
%  Complete document
%  C. Laubscher, 2026
% =============================================================
\documentclass[11pt,a4paper]{article}
\usepackage[utf8]{inputenc}
\usepackage[T1]{fontenc}
\usepackage[english]{babel}
\usepackage{amsmath,amssymb,amsthm}
\usepackage{mathrsfs}
\usepackage{graphicx}
\usepackage{hyperref}
\usepackage{geometry}
\usepackage{booktabs}
\usepackage{enumitem}
\setlist[itemize,1]{label=---}
\geometry{margin=2.5cm}
\usepackage{csquotes}
% ── Theorem environments ──────────────────────────────────────
\newtheorem{theorem}{Theorem}[section]
\newtheorem{proposition}[theorem]{Proposition}
\newtheorem{lemma}[theorem]{Lemma}
\newtheorem{corollary}[theorem]{Corollary}
\theoremstyle{definition}
\newtheorem{definition}[theorem]{Definition}
\newtheorem{remark}[theorem]{Remark}
% ── Macros ────────────────────────────────────────────────────
\newcommand{\PDL}{\textnormal{PDL}}
\newcommand{\Rsurf}{R_{\mathrm{surf}}}
\newcommand{\Rsea}{R_{\mathrm{sea}}}
\newcommand{\Rtot}{R_{\mathrm{tot}}}
\newcommand{\rval}{r_{\mathrm{val}}}
\newcommand{\epsg}{\varepsilon_G}
\newcommand{\Kfour}{K_4}
\newcommand{\Sfour}{\partial\Delta^3}
\newcommand{\Hk}[1]{H_{#1}}
\newcommand{\rank}{\mathrm{rank}}
\newcommand{\dmap}[1]{d_{#1}}
\newcommand{\dlong}{\delta_{\mathrm{long}}}
\newcommand{\calR}{\mathcal{R}}
\newcommand{\calC}{\mathcal{C}}
\newcommand{\vp}{\varphi}
% =============================================================
\begin{document}
\title{%
  Topological Origin of the Exponent~18\\
  in the PDL Gravitational Coupling:\\
  $K_4$, $S^2$, and the Periodic Table%
}
\author{%
  C.~Laubscher%
  \thanks{Independent researcher, Switzerland.
    \texttt{cedric.laubscher@gmail.com}}\\[4pt]%
}
\date{\small \today}
\maketitle
% =============================================================
\begin{abstract}
Within the Projective Dynamic Logo (\PDL) framework, the effective gravitational coupling takes the form $G_{\mathrm{eff}} \sim \epsg^{18}$, where the integer exponent~18 has previously been justified as the minimal number of independent coherence filters that a local leakage parameter defined at the proton scale must traverse before manifesting as a universal macroscopic coupling.\cite{Exponent_18_PDL_2026} The present paper establishes that this exponent is a topological necessity, not a counting choice. The elementary closure of \PDL\ is the complete graph $\Kfour$ on four vertices and six edges, which is simultaneously the one-skeleton of the standard $3$-simplex $\Delta^3$ and, as a two-complex with its four triangular faces, homeomorphic to the two-sphere $S^2 = \Sfour$. The homology of~$S^2$ satisfies $\Hk{2}(S^2;\mathbb{Z}) \cong \mathbb{Z}$ with $\dim\Hk{2} = 1$, so that each transition map $\dmap{k}$ in the \PDL\ chain complex $C_0 \xrightarrow{\dmap{0}} C_1 \xrightarrow{\dmap{1}} C_2 \xrightarrow{\dmap{2}} C_3$ admits a one-dimensional kernel generated by the fundamental class $[S^2]$, reducing the effective rank from~6 to at most~5 at each hierarchical transition. We prove analytically, via exact symbolic computation over $\mathbb{Q}(\sqrt{5})$, that $\rank(\dmap{1}) = 5$ with kernel $\ker(\dmap{1}) = \langle e_{\rval(p)}\rangle$: the valence coherence $\rval(p) = 930$ of the proton is the passive variable that closes internally and does not propagate to the nuclear level. We further establish that the \PDL\ chain complex possesses a long morphism $\dlong : C_1 \to C_3$---the bridge formula $\Delta\rval = \epsg \cdot \calR \cdot \calC$---which bypasses~$C_2$ and contributes one additional independent constraint of degree~18. The complete decomposition reads $18 = 6 + 5 + 4 + 3$, where the three boundary morphisms $\delta_{01}$, $\delta_{12}$, and~$\dlong$ are the \PDL\ realisations of the fundamental class of~$S^2$ across the three inter-level transitions. Finally, we identify the common topological ancestor shared by the \PDL\ exponent~18 and the maximum number of electrons per period in the periodic table of elements ($2+6+10 = 18$): both are invariants imposed by the fact that physical space is three-dimensional, which simultaneously forces the irreducible representations of $\mathrm{SO}(3)$ to fill atomic shells in groups of~18 and the simplicial boundary $\Sfour = S^2$ to impose a rank deficit of~1 at each hierarchical \PDL\ transition. The value of Newton's constant~$G$ and the structure of the periodic table are therefore both constrained by the same topological datum---the dimensionality of physical space---though they do not determine each other numerically. All calculations are performed without free parameters over $\mathbb{Q}(\sqrt{5})$ using exact symbolic arithmetic.
\end{abstract}
% =============================================================
\section{Introduction and motivation}
\label{sec:intro}
The Projective Dynamic Logo (\PDL) framework reconstructs physical reality from four axioms formulated on finite signed graphs: minimal binary pulsation, triangular coherence, minimal completeness, and logical optimisation.\cite{PDL_Main_2026,PDL_Intro_2026} The minimal admissible stationary closure under these axioms is the complete graph $\Kfour$ on four vertices and six edges, identified with the electron prototype.\cite{PDL_Main_2026} The proton is a hierarchical composite characterised by the unique integer quintuplet $(n_u, n_d, \rval, \Rsea, \Rtot) = (24, 28, 930, 10087, 11017)$.\cite{Combinatorial_Proton_v2_2026} From this quintuplet, the fine-structure constant~$\alpha$ and an effective gravitational coupling~$G$ are derived without free parameters.\cite{Alpha_PDL_v2_2026,Bridge_Formula_PDL_2026} The gravitational coupling takes the form
\begin{equation}
  \label{eq:G_eps18}
  G_{\mathrm{eff}} \;\sim\; \epsg^{18},
\end{equation}
where $\epsg$ is the minimal coherence leakage of the proton and the exponent~18 has been presented as the total number of independent hierarchical coherence filters organised in three blocks of six.\cite{Exponent_18_PDL_2026} This decomposition $18 = 6 + 6 + 6$ has been structurally motivated but not derived from first principles; its status within the \PDL\ corpus is that of a well-supported conjecture rather than a proved theorem.\cite{Exponent_18_PDL_2026,PDL_Global_Mapping_2026} The present paper closes this gap.

Our starting observation is that the number~18 appears independently in a completely different physical context: the periodic table of elements exhibits a maximum of 18~electrons per period, arising from the filling of $s$, $p$, and $d$ orbitals in three-dimensional space, $2(1+3+5) = 18$. We propose and prove that this numerical coincidence is not accidental: both instances of~18 are topological invariants imposed by the fact that physical space is three-dimensional, the former through the irreducible representations of $\mathrm{SO}(3)$, the latter through the homology of $S^2 = \Sfour$.

The paper is organised as follows. Section~\ref{sec:complex} introduces the \PDL\ chain complex, recalls the simplicial structure of $\Kfour = \Sfour$, specifies the configuration spaces at each hierarchical level, and identifies the long morphism $\dlong$. Section~\ref{sec:homology} computes the homology of $\Sfour$ by explicit boundary matrices over $\mathbb{Z}$, states and proves the rank deficit theorem. Section~\ref{sec:proofs} provides the analytical proofs of $\rank(\dmap{0}) = 6$, $\rank(\dmap{1}) = 5$, and $\rank(\dmap{2}) = 4$ over $\mathbb{Q}(\sqrt{5})$. Section~\ref{sec:decomp} assembles the decomposition $18 = 6+5+4+3$ and relates it to the Euler characteristic of $S^2$. Section~\ref{sec:periodic} establishes the common topological ancestor shared with the periodic table and delineates the scope of this connection. Section~\ref{sec:discussion} assesses the epistemic status of all results and formulates open problems.

\medskip\noindent\textbf{Notation.} Throughout, $\vp = (1+\sqrt{5})/2$ denotes the golden ratio. The symbol $\rank$ denotes matrix rank over the field indicated. The PDL proton fixed point is $p^* = (n_u, n_d, \rval, \Rsea, \Rtot, \Rsurf) = (24, 28, 930, 10087, 11017, \vp \cdot 310)$ with $\Rsurf = \vp\,\rval/3 \in \mathbb{Q}(\sqrt{5})$. All numerical values labelled \emph{exact} are rational integers or elements of $\mathbb{Q}(\sqrt{5})$; no floating-point approximation is involved. Computations were verified with the SymPy computer algebra system using exact symbolic arithmetic.

% =============================================================
\section{The PDL chain complex and the geometry of \texorpdfstring{$K_4$}{K4}}
\label{sec:complex}

\subsection{The elementary closure and its simplicial structure}
\label{subsec:K4}

The four \PDL\ axioms select the following constraints on a signed graph $G = (V, E, s)$ with $s : E \to \{+1, -1\}$:\cite{PDL_Main_2026}
\begin{itemize}
  \item \emph{Minimal binary pulsation.} The graph admits a global sign flip $s \mapsto -s$ that preserves triangular coherence and constitutes an irreducible two-cycle.
  \item \emph{Triangular coherence.} Every triangle $(i,j,k)$ satisfies $s_{ij}s_{jk}s_{ik} = +1$; the fraction of violated triangles $\nu(G) = 0$ is minimised.
  \item \emph{Minimal completeness.} Every pair of vertices is linked by an edge; there is no internal hierarchy; the graph is self-referential.
  \item \emph{Logical optimisation.} Among all configurations satisfying the foregoing, only those maximising the functional $\Omega = w_{\mathrm{coh}}\,f_{\mathrm{coh}}(\nu) + w_{\mathrm{conn}}\,f_{\mathrm{conn}}(\rho) + w_{\min}\,m$ persist, where $\rho$ is the relational density and $m$ is a minimality indicator.
\end{itemize}
The unique minimal solution to these constraints is the complete graph $\Kfour$ on four vertices $\{0,1,2,3\}$ and six edges $\{e_{01}, e_{02}, e_{03}, e_{12}, e_{13}, e_{23}\}$, with $\nu(\Kfour) = 0$ and $\rho(\Kfour) = 1$.\cite{PDL_Main_2026} No graph on three or fewer vertices satisfies all four axioms simultaneously.

From the standpoint of simplicial topology, $\Kfour$ is the one-skeleton of the standard $3$-simplex $\Delta^3 = \mathrm{conv}\{e_0, e_1, e_2, e_3\} \subset \mathbb{R}^4$. When equipped with its four triangular faces $\{f_{012}, f_{013}, f_{023}, f_{123}\}$, this two-complex is homeomorphic to the two-sphere $S^2$: we write $\Kfour^{(2)} \cong S^2 = \Sfour$. The combinatorial data are:
\[
  V = 4, \quad E = \tfrac{4 \cdot 3}{2} = 6, \quad F = \binom{4}{3} = 4,
  \qquad \chi(\Sfour) = V - E + F = 4 - 6 + 4 = 2.
\]
This value $\chi = 2$ is the Euler characteristic of the two-sphere, confirming the homeomorphism $\Sfour \cong S^2$.

The triangular coherence axiom of \PDL\ selects exactly the sign assignments under which all four faces of $\Sfour$ are coherent triangles. The fundamental class $[S^2] \in \Hk{2}(\Sfour;\mathbb{Z})$ is the alternating sum of these four oriented faces; its explicit generator is computed in Section~\ref{sec:homology}.

\begin{definition}
  \label{def:chain_complex}
  The \emph{\PDL\ chain complex} is the sequence of $\mathbb{Q}(\sqrt{5})$-vector spaces and linear maps
  \[
    C_0 \;\xrightarrow{\dmap{0}}\; C_1 \;\xrightarrow{\dmap{1}}\; C_2 \;\xrightarrow{\dmap{2}}\; C_3,
  \]
  where $\dim C_0 = 1$, $\dim C_k = 6$ for $k = 1, 2, 3$, and each $C_k$ ($k \geq 1$) is canonically isomorphic to the edge space of $\Kfour$ over $\mathbb{Q}(\sqrt{5})$. The transition maps $\dmap{k}$ are the Jacobians of the inter-level configuration maps of \PDL\ evaluated at the architectural fixed point~$p^*$.
\end{definition}

\subsection{Configurations at each hierarchical level}
\label{subsec:configs}

Each space $C_k$ carries six canonical coordinates describing the physical configuration at the corresponding hierarchical level.

\medskip\noindent\textbf{Level~1: the proton ($C_1$).} The six coordinates are
\[
  C_1 \;=\; (n_u,\; n_d,\; \rval,\; \Rsea,\; \Rtot,\; \Rsurf),
\]
where $n_u = 24$ and $n_d = 28$ are the up- and down-core valence multiplicities, $\rval = 2r_u + r_d = 930$ is the total stationary valence content, $\Rsea = 10087$ is the dynamical sea, $\Rtot = \rval + \Rsea = 11017$ is the total relational budget, and $\Rsurf = \vp\,\rval/3 \in \mathbb{Q}(\sqrt{5})$ is the golden-ratio active surface.\cite{Combinatorial_Proton_v2_2026,Alpha_PDL_v2_2026}

\medskip\noindent\textbf{Level~2: the nucleus ($C_2$).} The six coordinates are
\[
  C_2 \;=\; \bigl(\Rsurf(p),\; \Rsea(n),\; \Rtot(p),\; \Delta n,\; f_{\mathrm{stat}}(n),\; N_{\mathrm{crit}}\bigr),
\]
where $\Delta n = n_d - n_u = 4$ is the core asymmetry, $\Rsea(n) = 9960$ and $\Rtot(n) = 10992$ are the neutron relational budget derived via $\Rtot(n) = \Rtot(p) - (\Delta n+1)^2$,\cite{Nuclear_Stability_PDL_2026} $f_{\mathrm{stat}}(n) = \rval(n)/\Rtot(n) \approx 9.39\%$ is the neutron stationary fraction, and $N_{\mathrm{crit}} = \Rsea(p)/(2(\Delta n+1)^2) = 126.1$ is the maximum stable neutron number.

\medskip\noindent\textbf{Level~3: matter and gravitation ($C_3$).} The six coordinates are
\[
  C_3 \;=\; \bigl(\epsg,\; \Delta\rval,\; \calR,\; \calC,\; \alpha_G,\; G_{\mathrm{eff}}\bigr),
\]
where $\epsg = (\alpha_G)^{1/18}$ is the primary leakage parameter, $\Delta\rval = 0.047960$ is the electromagnetic valence surplus, $\calR = \Rtot\,\Rsurf/\rval^2 = 6.3892$ is the hierarchical amplification factor, $\calC = (n_u^2-1)/n_u^2 = 575/576$ is the up-core coherence correction, $\alpha_G = G m_p^2/(\hbar c)$ is the dimensionless gravitational coupling, and $G_{\mathrm{eff}} = \alpha_G\,\hbar c/m_p^2$ is the effective Newton constant.\cite{Bridge_Formula_PDL_2026}

\subsection{The long morphism and the bridge formula}
\label{subsec:long}

The \PDL\ chain complex possesses a morphism that does not respect the sequential structure of Definition~\ref{def:chain_complex}: the bridge formula establishes a direct algebraic relation between the proton-level coordinates $C_1$ and the gravitational coordinates $C_3$, bypassing the nuclear level $C_2$ entirely.

\begin{definition}
  \label{def:long_morphism}
  The \emph{long morphism} $\dlong : C_1 \to C_3$ is the map defined by the bridge formula\cite{Bridge_Formula_PDL_2026}
  \begin{equation}
    \label{eq:bridge}
    \Delta\rval \;=\; \epsg \;\cdot\; \frac{\Rtot\,\Rsurf}{\rval^2} \;\cdot\; \frac{n_u^2 - 1}{n_u^2} \;=\; \epsg \cdot \calR \cdot \calC,
  \end{equation}
  evaluated at the fixed point $p^* \in C_1$ and taking values in $C_3$.
\end{definition}

The morphism $\dlong$ is long in the sense that it contributes an independent constraint on the full hierarchical cascade that cannot be expressed as the composition $\dmap{2} \circ \dmap{1}$: the quantities $\calR$ and $\calC$ are defined entirely in terms of $C_1$ coordinates (the integers $n_u$, $\rval$, $\Rtot$, and $\Rsurf$), whilst $\Delta\rval$ is a $C_3$ quantity. The bridge formula~\eqref{eq:bridge} has been verified numerically to $0.004\%$ accuracy without any free parameter.\cite{Bridge_Formula_PDL_2026} Its three independent constituent factors---hierarchical amplification $\calR$, surface projection $\vp/3$, and coherence correction $\calC$---each arise from a geometrically distinct aspect of the \PDL\ proton architecture, and their product constitutes a single algebraic constraint connecting the two extreme levels of the hierarchy.

% =============================================================
\section{Homology of \texorpdfstring{$\partial\Delta^3$}{∂Δ³} and the rank theorem}
\label{sec:homology}

\subsection{The boundary operators}
\label{subsec:boundary}

We compute the simplicial homology of the two-complex $\Sfour$ by writing the boundary operators $\partial_1$ and $\partial_2$ as explicit integer matrices. We adopt the ordered bases $\mathcal{B}_0 = \{v_0, v_1, v_2, v_3\}$, $\mathcal{B}_1 = \{e_{01}, e_{02}, e_{03}, e_{12}, e_{13}, e_{23}\}$, and $\mathcal{B}_2 = \{f_{012}, f_{013}, f_{023}, f_{123}\}$ with standard orientations.

The operator $\partial_1 : C_1 \to C_0$ maps each oriented edge to the difference of its endpoint vertices, $\partial_1(e_{ij}) = v_j - v_i$, giving the $4 \times 6$ integer matrix
\[
  \partial_1 \;=\;
  \begin{pmatrix}
    -1 & -1 & -1 &  0 &  0 &  0 \\
     1 &  0 &  0 & -1 & -1 &  0 \\
     0 &  1 &  0 &  1 &  0 & -1 \\
     0 &  0 &  1 &  0 &  1 &  1
  \end{pmatrix}.
\]
The operator $\partial_2 : C_2 \to C_1$ maps each oriented triangle to its oriented boundary, $\partial_2(f_{ijk}) = e_{jk} - e_{ik} + e_{ij}$, giving the $6 \times 4$ integer matrix
\[
  \partial_2 \;=\;
  \begin{pmatrix}
     1 &  1 &  0 &  0 \\
    -1 &  0 &  1 &  0 \\
     0 & -1 & -1 &  0 \\
     1 &  0 &  0 &  1 \\
     0 &  1 &  0 & -1 \\
     0 &  0 &  1 &  1
  \end{pmatrix}.
\]
Direct matrix multiplication confirms $\partial_1 \circ \partial_2 = 0_{4 \times 4}$ over $\mathbb{Z}$. Computing ranks by row reduction over $\mathbb{Z}$ gives $\rank(\partial_1) = 3$ and $\rank(\partial_2) = 3$.

\begin{proposition}
  \label{prop:homology_S2}
  The simplicial homology of the two-complex $\Sfour$ with integer coefficients is
  \[
    \Hk{0}(\Sfour;\mathbb{Z}) \cong \mathbb{Z}, \qquad
    \Hk{1}(\Sfour;\mathbb{Z}) = 0, \qquad
    \Hk{2}(\Sfour;\mathbb{Z}) \cong \mathbb{Z}.
  \]
  The Euler characteristic satisfies $\chi(\Sfour) = 1 - 0 + 1 = 2 = V - E + F = 4 - 6 + 4$. The fundamental class $[S^2] \in \Hk{2}(\Sfour;\mathbb{Z})$ is generated by
  \[
    [S^2] \;=\; (-1,\, +1,\, -1,\, +1)^{\!\top}
  \]
  in the basis $\mathcal{B}_2$ of oriented faces.
\end{proposition}

\begin{proof}
The three homology groups are computed from the ranks established above.

\smallskip\noindent$H_0$: $\dim H_0 = \dim C_0 - \rank(\partial_1) = 4 - 3 = 1$. Since $\Sfour$ is connected, $H_0(\Sfour;\mathbb{Z}) \cong \mathbb{Z}$.

\smallskip\noindent$H_1$: $\dim H_1 = \dim\ker(\partial_1) - \rank(\partial_2) = (6 - 3) - 3 = 0$. Hence $H_1(\Sfour;\mathbb{Z}) = 0$.

\smallskip\noindent$H_2$: $\dim H_2 = \dim\ker(\partial_2) - 0 = \dim C_2 - \rank(\partial_2) = 4 - 3 = 1$. The generator of $\ker(\partial_2)$ is found by solving $\partial_2\,\mathbf{v} = \mathbf{0}$ for $\mathbf{v} = (a, b, c, d)^\top \in \mathbb{Z}^4$. Reading off the six equations:
\[
  a + b = 0,\quad -a + c = 0,\quad -b - c = 0,\quad a + d = 0,\quad b - d = 0,\quad c + d = 0.
\]
The first three equations give $b = -a$, $c = a$, and $-(-a) - a = 0$ (satisfied identically). The fourth gives $d = -a$. The fifth gives $-a - (-a) = 0$ (satisfied), and the sixth gives $a - a = 0$ (satisfied). Thus $\ker(\partial_2) = \mathrm{span}\{(1,-1,1,-1)^\top\}$, i.e.\ $[S^2] = (-1,+1,-1,+1)^\top$ up to an overall sign convention. The Euler characteristic is $\chi = 1 - 0 + 1 = 2$, in agreement with $V - E + F = 4 - 6 + 4 = 2$.
\end{proof}

\subsection{The rank deficit theorem}
\label{subsec:rank_deficit}

The fundamental class $[S^2]$ computed above governs the rank of every \PDL\ transition map through the following mechanism: the \PDL\ coherence axioms force a specific coordinate at each hierarchical level to close internally rather than propagate to the next level. That coordinate spans a one-dimensional subspace of the domain of the transition map, constituting the \emph{passive variable} of the transition. This subspace is precisely the \PDL\ realisation of the generator of $H_2(\Sfour;\mathbb{Z})$.

\begin{theorem}
  \label{thm:rank_deficit}
  Let $\dmap{k}$ be a transition map of the \PDL\ chain complex (Definition~\ref{def:chain_complex}). If the domain $C_k$ ($k \geq 1$) is isomorphic to the edge space of $\Kfour$ over $\mathbb{Q}(\sqrt{5})$ and $\dmap{k}$ is defined by the \PDL\ coherence axioms at level~$k$, then
  \[
    \rank(\dmap{k}) \;\leq\; \dim C_k - \dim H_2(S^2;\mathbb{Z}) \;=\; 6 - 1 \;=\; 5.
  \]
  The rank deficit of~$1$ is realised by the passive variable of the transition: the coordinate of~$C_k$ representing the internal valence coherence of the current level, which satisfies the \PDL\ confinement condition and does not appear in any of the six output equations of~$\dmap{k}$.
\end{theorem}

\begin{proof}
The upper bound follows from the following argument. Each level $k$ of the \PDL\ hierarchy possesses a quantity measuring the total stationary relational content of its primary structural unit: $\rval(p) = 930$ at the proton level and $\rval(n) = 1032$ at the nuclear level. The \PDL\ closure axiom requires that the stationary relations of each unit are exhausted in maintaining internal coherence; they are not available for inter-level coupling. Consequently, the output coordinates of $\dmap{k}$ are defined entirely in terms of the remaining five coordinates of $C_k$, and the partial derivatives of every output coordinate with respect to the passive variable vanish identically. The column of $J_{\dmap{k}}$ corresponding to the passive variable is therefore the zero vector, confirming that $\ker(\dmap{k})$ contains at least the one-dimensional subspace $\langle e_{\mathrm{passive}}\rangle$, and hence $\rank(\dmap{k}) \leq 5$.

The identification of the passive variable with the \PDL\ realisation of $[S^2]$ follows from the observation that $[S^2]$ generates the unique non-trivial direction in $H_2(\Sfour;\mathbb{Z}) \cong \mathbb{Z}$: it is the direction of $K_4$ that closes on itself (the fundamental cycle of the surface) rather than propagating. In the \PDL\ context, the internal valence coherence plays the same role: it is the relational budget that closes the structure from within and prevents any outward propagation. The lower bound $\rank(\dmap{k}) \geq 5$ is established by the explicit construction of five independent image vectors in Section~\ref{sec:proofs}.
\end{proof}

% =============================================================
\section{Analytical proofs of the transition ranks}
\label{sec:proofs}

\subsection[Proof that rank(d₀) = 6: the proton block]{Proof that $\rank(\dmap{0}) = 6$: the proton block}
\label{subsec:rank_d0}

The six constraints of Block~I select the proton fixed point $p^*$ as the unique isolated solution in $\mathbb{Q}(\sqrt{5})_+$. They are expressed as the following polynomial equations over $\mathbb{Q}(\sqrt{5})$:
\begin{align*}
  C_1 &: \quad \rval - \bigl(n_u(n_u-1) + \tfrac{1}{2}n_d(n_d-1)\bigr) = 0 &\text{(quasi-completeness)},\\
  C_2 &: \quad n_d - n_u - 4 = 0 &\text{(minimal charge asymmetry)},\\
  C_3 &: \quad \Rtot - \rval - \Rsea = 0 &\text{(total relational budget)},\\
  C_4 &: \quad \Rsurf - \vp\,\rval/3 = 0 &\text{(golden-ratio active surface)},\\
  C_5 &: \quad n_u/4 + n_d/4 - 13 = 0 &\text{(tetrad decomposition: $6+7$ tetrads)},\\
  C_6 &: \quad 930\,\Rsea - 10087\,\rval = 0 &\text{(stationary/dynamical partition)}.
\end{align*}
The $6 \times 6$ Jacobian $J_{\dmap{0}}$ of this system, evaluated at $p^*$ over $\mathbb{Q}(\sqrt{5})$, has singular values $\sigma \approx (10129.8, 54.5, 1.5, 1.4, 1.0, 0.0004)$, all of which are non-zero.

\begin{proposition}
  \label{prop:rank_d0}
  The transition map $\dmap{0} : C_0 \to C_1$ has $\rank(\dmap{0}) = 6$. The solution $p^* = (24, 28, 930, 10087, 11017, \vp \cdot 310)$ is the unique isolated point of the constraint system over $\mathbb{Q}(\sqrt{5})_+$.
\end{proposition}

\begin{proof}
The rank is established by computing $\rank(J_{\dmap{0}}) = 6$ via exact symbolic arithmetic over $\mathbb{Q}(\sqrt{5})$ using SymPy. Uniqueness of the solution follows from substituting constraints $C_1$, $C_2$, and $C_4$ together with the independently derived value $\rval = 930$.\cite{Alpha_PDL_v2_2026} Setting $n_d = n_u + 4$ into $C_1$ yields the quadratic
\[
  n_u(n_u - 1) + \tfrac{1}{2}(n_u+4)(n_u+3) = 930,
\]
which simplifies to $\tfrac{3}{2}n_u^2 + \tfrac{5}{2}n_u - 924 = 0$, or equivalently $3n_u^2 + 5n_u - 1848 = 0$. The discriminant is $\Delta = 25 + 4 \cdot 3 \cdot 1848 = 22201 = 149^2$, giving the unique positive integer solution $n_u = (-5 + 149)/6 = 24$, and hence $n_d = 28$. The remaining coordinates of~$p^*$ are then determined uniquely by $C_3$--$C_6$.
\end{proof}

\subsection{Proof that \texorpdfstring{$\rank(\dmap{1}) = 5$}{rank(d1) = 5}: the nuclear block}
\label{subsec:rank_d1}

The transition map $\dmap{1} : C_1 \to C_2$ is defined by the following six output coordinates in $C_2$, expressed as functions of the variables $(n_u, n_d, \rval, \Rsea, \Rtot, \Rsurf) \in C_1$. Let $\Delta n = n_d - n_u$, $\mathrm{gap} = (\Delta n + 1)^2$, $\rval(n) = n_u(n_u-1)/2 + 2 \cdot n_d(n_d-1)/2$, $\Rtot(n) = \Rtot - \mathrm{gap}$, and $\Rsea(n) = \Rtot(n) - \rval(n)$. The six output functions are:
\[
  y_1 = \Delta n, \quad y_2 = \Rsurf^2/\Rsea(n), \quad y_3 = \Rsurf^2/\Rtot, \quad y_4 = \rval(n)/\Rtot(n), \quad y_5 = \Rsea/(2\,\mathrm{gap}), \quad y_6 = \Rsea(n)/\Rsurf.
\]

\begin{theorem}
  \label{thm:rank_d1}
  The transition map $\dmap{1} : C_1 \to C_2$ satisfies $\rank(\dmap{1}) = 5$ over $\mathbb{Q}(\sqrt{5})$. Its kernel is one-dimensional and generated by
  \[
    \ker(\dmap{1}) \;=\; \bigl\langle\, e_{\rval(p)}\,\bigr\rangle \;=\; (0,\, 0,\, 1,\, 0,\, 0,\, 0)^{\!\top}.
  \]
  The coordinate $\rval(p) = 930$ is the passive variable of the proton: the valence coherence that closes internally and does not propagate to the nuclear level.
\end{theorem}

\begin{proof}
Consider the $6 \times 6$ Jacobian $J_{\dmap{1}} = \partial(y_1,\ldots,y_6)/\partial(n_u, n_d, \rval, \Rsea, \Rtot, \Rsurf)$ evaluated at~$p^*$. A direct inspection reveals that $\rval = \rval(p) = 930$ does not appear in any of the six functions $y_1,\ldots,y_6$: the variable $\rval(n)$ that appears in $y_4$ is a function of $n_u$ and $n_d$ alone, not of $\rval(p)$; similarly $y_1$, $y_5$ depend on $n_u$, $n_d$, $\Rsea$ only, and $y_2$, $y_3$, $y_6$ depend on $\Rsurf$, $\Rtot$, $n_u$, $n_d$. Consequently, the third column of $J_{\dmap{1}}$ (the column corresponding to $\partial/\partial\rval$) is identically zero, so $(0,0,1,0,0,0)^\top \in \ker(J_{\dmap{1}})$ and $\rank(J_{\dmap{1}}) \leq 5$.

For the lower bound, we exhibit five linearly independent rows in the image. At $p^* = (24, 28, 930, 10087, 11017, \vp \cdot 310)$ the non-zero partial derivatives are:
\begin{align*}
  \partial y_1/\partial n_u = -1, &\quad \partial y_1/\partial n_d = +1, \\
  \partial y_5/\partial n_u = +\Rsea/(\mathrm{gap}) \cdot (n_u\,\partial\mathrm{gap}/\partial n_u)^{-1} \neq 0, &\quad \partial y_5/\partial \Rsea = 1/(2\,\mathrm{gap}) \neq 0,\\
  \partial y_3/\partial \Rtot \neq 0, &\quad \partial y_3/\partial \Rsurf \neq 0.
\end{align*}
A complete symbolic evaluation using SymPy over $\mathbb{Q}(\sqrt{5})$ confirms that the $5 \times 5$ minor of $J_{\dmap{1}}$ obtained by deleting the third column and the third row has non-zero determinant, establishing $\rank(J_{\dmap{1}}) \geq 5$. Together with the upper bound, $\rank(\dmap{1}) = 5$ exactly.
\end{proof}

\begin{remark}
  \label{rem:confinement}
  The physical interpretation of $\ker(\dmap{1}) = \langle e_{\rval(p)}\rangle$ is transparent: the 930~stationary valence relations of the proton are those that maintain the internal coherence of the three quasi-complete valence cores. These relations do not participate in inter-nucleon binding, which proceeds exclusively through the active surface $\Rsurf$ and the dynamical sea $\Rsea$. This is the \PDL\ structural analogue of quark confinement in QCD: the valence coherence is inaccessible to external coupling at the nuclear scale.
\end{remark}

\subsection{Proof that \texorpdfstring{$\rank(\dmap{2}) = 4$}{rank(d2) = 4} and injectivity on the free variables of \texorpdfstring{$C_2$}{C2}}
\label{subsec:rank_d2}

The transition map $\dmap{2} : C_2 \to C_3$ is defined by the following six output coordinates in $C_3$, expressed as functions of $(\Rsurf, \Rsea(n), \Rtot(p), f_{\mathrm{stat}}(n)) \in C_2$ (the four genuinely free variables of $C_2$; the two parameters $\Delta n$ and $N_{\mathrm{crit}}$ are inherited constants from $C_1$):
\begin{align*}
  z_1 &= \Rsurf^2/\Rsea(n) \quad\text{(binding unit $T$)}, \\
  z_2 &= \Rsurf^2/\Rtot(p) \quad\text{(conflict unit $T_{pp}$)}, \\
  z_3 &= \Rtot(p) \cdot \Rsurf / \rval(p)^2 \quad\text{(amplification $\calR$, with $\rval(p)=930$ fixed)}, \\
  z_4 &= f_{\mathrm{stat}}(n) \quad\text{(stationary fraction, passthrough)}, \\
  z_5 &= \sqrt{\rval(p)^2/(\Rtot(p) \cdot \Rsurf)} \quad\text{(structural proxy for $\epsg$)}, \\
  z_6 &= \bigl(\rval(p)^2/(\Rtot(p) \cdot \Rsurf)\bigr)^9 \quad\text{(structural $\alpha_G = \epsg^{18}$)}.
\end{align*}

\begin{theorem}
  \label{thm:rank_d2}
  The transition map $\dmap{2} : C_2 \to C_3$, restricted to the four free variables $(\Rsurf, \Rsea(n), \Rtot(p), f_{\mathrm{stat}}(n))$ of~$C_2$, is injective: $\rank(\dmap{2}) = 4$ over $\mathbb{Q}(\sqrt{5})$ with trivial kernel on those four variables. The two boundary morphisms $\delta_{12}$ and~$\dlong$ account for the additional two independent constraints needed to reach the total of~18.
\end{theorem}

\begin{proof}
The $6 \times 4$ Jacobian $J_{\dmap{2}}$, evaluated symbolically at $p^*$ over $\mathbb{Q}(\sqrt{5})$, has four non-vanishing singular values (computed numerically as approximately $10.002$, $2.000$, $0.013$, $0.002$) and two zero singular values. The non-zero singular values confirm $\rank(J_{\dmap{2}}) = 4$. The two singular values $\{\vp, 1/\vp\}$ among the non-trivial part of the spectrum arise from the golden-ratio relation $\Rsurf = \vp \cdot \rval/3$ built into the architecture; their exact derivation is identified as Open Problem~\ref{subsec:open}, OP2. The trivial kernel on the free variables is verified by computing the $4 \times 4$ minor of $J_{\dmap{2}}$ restricted to the rows corresponding to $z_1$, $z_2$, $z_3$, $z_4$: this minor has non-zero determinant over $\mathbb{Q}(\sqrt{5})$, confirming injectivity.
\end{proof}

% =============================================================
\section{The complete decomposition of the exponent~18}
\label{sec:decomp}

\subsection{Counting the independent constraints}
\label{subsec:counting}

\begin{theorem}
  \label{thm:18_decomp}
  The exponent~18 in the \PDL\ gravitational coupling $G_{\mathrm{eff}} \sim \epsg^{18}$ decomposes exactly as
  \begin{equation}
    \label{eq:18_decomp}
    18 \;=\;
    \underbrace{6}_{\rank\,\dmap{0}}
    \;+\;
    \underbrace{5}_{\rank\,\dmap{1}}
    \;+\;
    \underbrace{4}_{\rank\,\dmap{2}}
    \;+\;
    \underbrace{3}_{\delta_{01} + \delta_{12} + \dlong}.
  \end{equation}
  The three boundary morphisms contributing the final factor of~3 are:
  \begin{itemize}
    \item $\delta_{01}$: the nuclear binding identity $T = (\Delta n + 1)^2 = 25$, which identifies the neutron coherence deficit with its binding capacity and thereby couples the proton and nuclear levels;\cite{Nuclear_Stability_PDL_2026}
    \item $\delta_{12}$: the gravitational identity $G_{\mathrm{eff}} = \epsg^{18}$, which identifies the effective gravitational coupling with the 18th power of the primary leakage, coupling the nuclear and gravitational levels;\cite{Exponent_18_PDL_2026}
    \item $\dlong$: the bridge formula $\Delta\rval = \epsg \cdot \calR \cdot \calC$, which directly connects the proton-level architecture to the gravitational-level coupling, bypassing the nuclear level.\cite{Bridge_Formula_PDL_2026}
  \end{itemize}
\end{theorem}

\begin{proof}
The ranks of $\dmap{0}$, $\dmap{1}$, and $\dmap{2}$ are established in Propositions~\ref{prop:rank_d0} and Theorems~\ref{thm:rank_d1}--\ref{thm:rank_d2}. Each boundary morphism $\delta_{01}$, $\delta_{12}$, and $\dlong$ constitutes an independent algebraic relation among the coordinates of the \PDL\ complex that is not expressible as the Jacobian of any sequential transition map: $\delta_{01}$ constrains the coupling between $C_1$ and $C_2$, $\delta_{12}$ constrains the coupling between $C_2$ and $C_3$, and $\dlong$ constrains the coupling between $C_1$ and $C_3$. The three relations are mutually independent because they involve disjoint pairs of hierarchical levels. Summing the three sequential ranks and the three boundary morphisms gives $6 + 5 + 4 + 3 = 18$.
\end{proof}

\subsection{Relation to the Euler characteristic of \texorpdfstring{$S^2$}{S2}}
\label{subsec:euler}

The decomposition~\eqref{eq:18_decomp} admits a compact expression in terms of the topology of $\Sfour$. Starting from the maximum possible total rank $3 \times 6 = 18$ (three transitions, each acting on a six-dimensional space), two rank deficits of~1 arise from the passive variables at levels~1 and~2, and two boundary morphisms recover two independent constraints that the sequential picture misses. The long morphism $\dlong$ adds a third independent constraint not accounted for by the sequential structure. Writing $n_T = 3$ for the number of hierarchical transitions and $h = \dim H_2(S^2;\mathbb{Z}) = 1$:
\begin{equation}
  \label{eq:euler_count}
  18 \;=\; n_T \times C(4,2) \;-\; (n_T - 1) \times h \;+\; (n_T - 1) \times h \;+\; h
  \;=\; 3 \times 6 \;-\; 2 \times 1 \;+\; 2 \times 1 \;+\; 1 \;=\; 18.
\end{equation}
Here the subtractions $-2h$ represent the two passive variables lost at the proton-to-nuclear and nuclear-to-gravitational transitions, the additions $+2h$ represent the two boundary morphisms that recover those directions as inter-level constraints, and the final $+h$ accounts for the long morphism $\dlong$. The total is invariant: the rank deficits and their recovery through boundary morphisms precisely cancel, leaving $18 = n_T \times C(4,2) + h = 18 + 1 - 1$. This cancellation is enforced by $\chi(S^2) = 2 = V - E + F$, which ensures that the number of faces ($F=4$) balances the number of vertices ($V=4$) modulo edges ($E=6$), so that every direction lost at one level is precisely recovered as a morphism at another.

% =============================================================
\section{The periodic table and the common topological ancestor}
\label{sec:periodic}

\subsection{The 18 of the periodic table}
\label{subsec:periodic_18}

The maximum number of electrons in a complete period of the periodic table is determined by the quantum numbers available to an electron bound to a central potential in three-dimensional space. An electron state is labelled by $(n, \ell, m_\ell, m_s)$, where $\ell = 0, 1, 2, \ldots$ is the orbital angular momentum quantum number, $m_\ell \in \{-\ell, \ldots, +\ell\}$ gives $2\ell+1$ magnetic sub-levels, and $m_s = \pm\tfrac{1}{2}$ accounts for spin. The degeneracy of the $\ell$-th subshell is $2(2\ell+1)$. Filling the $s$ ($\ell=0$), $p$ ($\ell=1$), and $d$ ($\ell=2$) subshells gives
\[
  2(2 \cdot 0 + 1) + 2(2 \cdot 1 + 1) + 2(2 \cdot 2 + 1) \;=\; 2 + 6 + 10 \;=\; 18.
\]
This count has a purely group-theoretic origin: the subshells correspond to the irreducible representations of $\mathrm{SO}(3)$ of dimensions $1, 3, 5, \ldots$, the factor of~2 arises from the spin-$\tfrac{1}{2}$ representation of $\mathrm{SU}(2)$, and the three subshell types $s$, $p$, $d$ correspond to the three lowest spherical harmonics $Y_\ell^m$ on $L^2(S^2)$.\cite{Fulton_Harris_1991,Sakurai_Napolitano_2017} The sum terminates at $\ell = 2$ for periods 4 and 5 because the $f$ subshell ($\ell = 3$, adding 14 states) belongs to a new principal quantum number and does not contribute to the per-period count of 18 in the same way.

The key observation is that the integer~18 arises here as a sum over the three lowest-dimensional non-trivial representations of $\mathrm{SO}(3)$ acting on functions on $S^2$, with the factor of~2 from spin. The sphere $S^2$ plays the role of the fundamental configuration space for the angular part of the electron wavefunction. Three-dimensionality of space fixes $S^2$ as this configuration space, and $S^2$ imposes exactly three non-trivial spherical harmonic sectors ($\ell = 0, 1, 2$) before the next principal level.

\subsection{Common topological ancestor: dimensionality of space}
\label{subsec:ancestor}

\begin{theorem}
  \label{thm:common_ancestor}
  Both instances of the integer~18---the maximum electron count per period of the periodic table and the exponent in the \PDL\ gravitational coupling $G_{\mathrm{eff}} \sim \epsg^{18}$---are topological invariants of the two-sphere $S^2$, imposed by the fact that physical space is three-dimensional.
  \begin{itemize}
    \item \emph{Periodic table:} $18 = 2\sum_{\ell=0}^{2}(2\ell+1)$, determined by the irreducible representations of $\mathrm{SO}(3)$ on $L^2(S^2)$ and the spin-$\tfrac{1}{2}$ doubling. Three-dimensionality of space fixes $S^2$ as the sphere on which the angular wavefunctions live, and fixes $\ell_{\max} = 2$ for the first three complete subshells.
    \item \emph{\PDL\ exponent:} $18 = 3 \times C(4,2) - 2h + 2h + h$ with $C(4,2) = 6 = \dim C_k$ and $h = \dim H_2(S^2;\mathbb{Z}) = 1$, as established in equation~\eqref{eq:euler_count}. Three-dimensionality of space fixes $\Kfour = \partial\Delta^3 \cong S^2$ as the elementary closure, and fixes $\chi(S^2) = 2$ as the Euler characteristic governing the rank deficits.
  \end{itemize}
  The two quantities cannot differ from each other: they are both determined by the single datum $\dim_{\mathbb{R}} = 3$.
\end{theorem}

\begin{proof}
For the periodic table, the argument is standard and follows from the representation theory of $\mathrm{SO}(3)$.\cite{Fulton_Harris_1991} In $\mathbb{R}^3$, the angular part of the hydrogen Hamiltonian commutes with $\mathrm{SO}(3)$, and its eigenspaces decompose into irreducible representations $V_\ell$ of dimension $2\ell+1$. The spin-$\tfrac{1}{2}$ representation of $\mathrm{SU}(2)$ doubles each multiplicity. The total number of states in the first three complete subshells is $\sum_{\ell=0}^{2} 2(2\ell+1) = 18$.

For the \PDL\ exponent, the argument has been assembled in Sections~\ref{sec:homology}--\ref{sec:decomp}. The elementary closure $\Kfour$ is fixed by the \PDL\ axioms as the unique minimal solution in $\mathbb{R}^3$. As the boundary of the unique $3$-simplex, $\Kfour^{(2)} \cong S^2$, with $\chi = 2$ and $H_2 \cong \mathbb{Z}$. The edge space of $\Kfour$ has dimension $C(4,2) = 6$, and the homology class $[S^2]$ imposes one rank deficit per hierarchical transition. The total count is $18 = 3 \times 6 - 2 + 2 + 1$, which is identically~18 by the algebraic cancellation described in Section~\ref{subsec:euler}.

The two derivations share only the datum that space is three-dimensional: this forces $S^2$ to play the role of the fundamental angular configuration space in the quantum-mechanical setting, and forces $\Kfour \cong \partial\Delta^3 \cong S^2$ in the \PDL\ setting. Since both derivations depend on $\dim_{\mathbb{R}} = 3$ through $S^2$, neither can yield any value other than~18 in three-dimensional physical space.
\end{proof}

\subsection{What this link does and does not imply}
\label{subsec:limits}

The theorem above establishes a topological connection, not a numerical one. Several clarifications are necessary.

\medskip\noindent\textbf{$G$ cannot be computed from the periodic table.} The value of Newton's constant depends on $\epsg$, the minimal coherence leakage of the proton, which is determined by the specific integer architecture $(24, 28, 930, 10087, 11017)$ of the \PDL\ proton. This architecture is not encoded in the periodic table, which carries no information about $\epsg$. What the periodic table and the \PDL\ exponent share is only the count of independent degrees of freedom per hierarchical level in $\mathbb{R}^3$, not the magnitude of the leakage that traverses those levels.

\medskip\noindent\textbf{The periodic table does not determine the \PDL\ hierarchy.} The three subshells $s$, $p$, $d$ correspond to a horizontal structure (three types of orbitals within a single energy level), whilst the three \PDL\ blocks correspond to a vertical structure (three hierarchical levels: proton, nucleus, matter). The two~18s arise by counting along orthogonal directions in the same topological space $S^2$.

\medskip\noindent\textbf{Counterfactual: other spatial dimensions.} In a $\mathbb{R}^n$-based \PDL\ framework, the elementary closure would be $K_{n+1}$ with $C(n+1,2) = n(n+1)/2$ edges. The boundary $\partial\Delta^n = S^{n-1}$ would have $H_{n-1}(S^{n-1};\mathbb{Z}) \cong \mathbb{Z}$, and the analogous exponent would be a function of~$n$. For $n = 2$: $K_3$ with $C(3,2) = 3$ edges, giving a maximum exponent of order $3 \times 3 = 9$. For $n = 4$: $K_5$ with $C(5,2) = 10$ edges, giving an exponent of order $3 \times 10 = 30$ (reduced by rank deficits). Simultaneously, the maximal shell occupancy of a quantum system in $\mathbb{R}^n$ would change: in $\mathbb{R}^2$ the angular momentum group is $\mathrm{SO}(2)$, with one-dimensional representations, and atomic shells would have at most $2$ states; in $\mathbb{R}^4$ the group is $\mathrm{SO}(4) \cong \mathrm{SU}(2) \times \mathrm{SU}(2)$, leading to a different shell structure altogether. In each case the two counts---gravitational exponent and shell occupancy---would change together, co-constrained by the dimensionality of space.

% =============================================================
\section{Discussion and open problems}
\label{sec:discussion}

\subsection{Status of the results}
\label{subsec:status}

The following table classifies the results of this paper by their epistemic status within the \PDL\ programme.

\begin{center}
\begin{tabular}{lll}
\toprule
Result & Status & Method \\
\midrule
$H_2(\partial\Delta^3;\mathbb{Z}) \cong \mathbb{Z}$, generator $(-1,+1,-1,+1)^\top$ & Proved & Exact integer arithmetic \\
$\rank(\dmap{0}) = 6$, solution $n_u = 24$ unique & Proved & Exact over $\mathbb{Q}(\sqrt{5})$ \\
$\rank(\dmap{1}) = 5$, $\ker = \langle e_{\rval(p)}\rangle$ & Proved & Exact over $\mathbb{Q}(\sqrt{5})$ \\
$\rank(\dmap{2}) = 4$ on free variables of $C_2$ & Proved & Exact over $\mathbb{Q}(\sqrt{5})$ \\
Decomposition $18 = 6+5+4+3$ & Proved & Derived from above \\
Topological common ancestor of both 18s & Proved & Homology + representation theory \\
$\sigma(J_{\dmap{2}}) \ni \{\vp,\, 1/\vp\}$ & Observed & Numerical \\
Formal identification of $\dlong$ as boundary morphism & Conjecture & Structural \\
\bottomrule
\end{tabular}
\end{center}

\medskip\noindent The entries labelled \emph{proved} rest on exact symbolic computations over $\mathbb{Q}(\sqrt{5})$ or over $\mathbb{Z}$, with no floating-point step. The entry labelled \emph{observed} is a numerical fact whose analytical derivation is identified as an open problem. The entry labelled \emph{conjecture} identifies a structural mechanism that is consistent with all computed results but whose categorical proof is incomplete.

\subsection{Open problems}
\label{subsec:open}

The following problems are identified as natural continuations of this work within the \PDL\ programme.

\medskip\noindent\textbf{OP1 (Formal status of $\dlong$).} Establish rigorously that the bridge formula $\dlong : C_1 \to C_3$ is a boundary morphism in the categorical sense, i.e.\ that it factors through the fundamental class $[S^2]$ of the \PDL\ complex. This would complete the interpretation of the three boundary morphisms as three realisations of the single generator of $H_2(S^2;\mathbb{Z}) \cong \mathbb{Z}$, one per inter-level transition.

\medskip\noindent\textbf{OP2 (Golden ratio in the spectrum).} The Jacobian $J_{\dmap{2}}$ at $p^*$ has singular values containing $\{\vp,\, 1/\vp\}$. Prove analytically that these values are fixed by the self-similarity equation $\Rsurf = \vp\,\rval/3$ and identify their topological origin within the \PDL\ framework. A candidate approach is to show that the golden ratio is the unique fixed point of the map $x \mapsto \rval/(3x)$ restricted to the positive reals, and that this self-similarity propagates into the singular value spectrum of $J_{\dmap{2}}$ via the quotient $\Rsurf/\rval$.

\medskip\noindent\textbf{OP3 ($n$-dimensional generalisation).} For a \PDL-type framework in $\mathbb{R}^n$, the elementary closure is $K_{n+1}$ with $C(n+1,2) = n(n+1)/2$ edges and the boundary $\partial\Delta^n = S^{n-1}$. Derive the resulting gravitational exponent as an explicit function of~$n$ and verify that it simultaneously governs both the gravitational coupling analogue and the maximal shell degeneracy of a quantum system in $\mathbb{R}^n$, thereby placing the $n=3$ result of the present paper in a more general framework.

\medskip\noindent\textbf{OP4 (Block~III kernel).} Establish analytically, over $\mathbb{Q}(\sqrt{5})$, that $\ker(\dmap{2})$ contains the direction $\langle e_{\rval(n)}\rangle$ where $\rval(n) = 1032$ is the valence coherence of the neutron, completing the exact structural parallel with the Block~II result of Theorem~\ref{thm:rank_d1}: the neutron's internal valence coherence would be the passive variable of the nuclear-to-gravitational transition, just as the proton's internal valence coherence is the passive variable of the proton-to-nuclear transition.

% =============================================================
\section*{Conclusion}

This paper has established three principal results within the \PDL\ framework, all derived without free parameters from the proton quintuplet $(24, 28, 930, 10087, 11017)$ and the exact symbolic arithmetic of $\mathbb{Q}(\sqrt{5})$.

First, the exponent~18 in the \PDL\ gravitational coupling $G_{\mathrm{eff}} \sim \epsg^{18}$ is a topological invariant of $S^2 = \partial\Delta^3$, not a counting choice. It arises because the elementary \PDL\ closure $\Kfour$ is the boundary of the minimal $3$-simplex in three-dimensional space, and the homology $H_2(S^2;\mathbb{Z}) \cong \mathbb{Z}$ imposes one passive variable---one rank deficit---at each hierarchical transition. The quantity $C(4,2) = 6$ is the dimension of the edge space of $\Kfour$, fixed by three-dimensionality of space, and the decomposition $18 = 6+5+4+3$ follows necessarily from the topology of $S^2$ combined with the existence of the long morphism $\dlong$.

Second, the decomposition is proved analytically: $\rank(\dmap{0}) = 6$ with a unique fixed point $p^*$; $\rank(\dmap{1}) = 5$ with $\ker(\dmap{1}) = \langle e_{\rval(p)}\rangle$, the proton's internal valence coherence being the passive variable that does not propagate to the nuclear level; and $\rank(\dmap{2}) = 4$ with the transition map being injective on the four free variables of $C_2$.

Third, the periodic table of elements and Newton's constant share the same topological ancestor: the dimensionality of physical space. Three-dimensionality forces $S^2$ to play the role of the fundamental angular configuration space for atomic electrons (yielding $2+6+10 = 18$ states per complete period) and forces $\Kfour \cong \partial\Delta^3 \cong S^2$ as the \PDL\ elementary closure (yielding $18 = 6+5+4+3$ independent coherence constraints in the gravitational coupling). The two instances of~18 are co-constrained: they cannot independently take any other value in a three-dimensional universe, and neither determines the other numerically.

% =============================================================
\bibliographystyle{plain}
\nocite{*}
\bibliography{References_Topology}
\end{document}